% ===================================================================
%  HOTO Na 16 — Babban shago. --- Kasuwar biki. Kayan wasa.
%
%  Traduction, faite en 2026, en haoussa d'aujourd'hui, sur l'ido de
%  texto/io. Regles de la colonne : voir 10-hoto-01.tex. Une seule
%  numerotation. Dernier tableau du livret.
%
%  ECART DECLARE, ET IL FAUT LE SAVOIR AVANT D'ECRIRE. Le bloc
%  t16-15-1 est l'un des trois que renvoji.py laisse passer sans rien
%  dire : le fac-simile ido saute la parenthese ouvrante de la lettre
%  « (d) » des hemispheres, que la planche grave pourtant. Un ecart
%  declare rend le bloc AVEUGLE. L'ordre 85 a b c d e f 86 87 88 89 a
%  donc ete pose a la main et recompte a la main, comme au tableau 8
%  pour le gobelet.
%
%  QUINZE DES VINGT-TROIS BETES DE LA MENAGERIE EN CARTON PORTENT UN
%  NOM HAOUSSA ORDINAIRE, ET C'EST LE PLUS FORT SCORE DU RELEVE. La
%  menagerie javanaise etait a moitie javanaise, la persane a moitie
%  iranienne ; celle-ci est haoussa aux trois quarts :
%
%    aku (6)       le perroquet   karkanda (7)  le rhinoceros
%    jimina (9)    l'autruche     biri (10)     le singe
%    dila (11)     le renard      kada (12)     le crocodile
%    raƙumi (23)   le dromadaire  damisa (18)   le leopard
%    kura (17)     l'hyene        kerkeci (21)  le loup
%    kusu (19)     le rat         ɓera (20)     la souris
%    dorina (22)   l'hippopotame  jemage (35)   la chauve-souris
%    mesa (36)     le boa         kunkuru (37)  la tortue
%
%  Et les huit qui manquent viennent TOUTES de l'eau froide ou du pays
%  froid : le renne, le phoque, la baleine, le requin, la belette, le
%  levrier. Le Sahel est presque entier sur cette planche, et ce qui
%  n'y est pas n'y a jamais mis les pieds.
%
%  ET LE DROMADAIRE (23) EST « raƙumi » TOUT COURT, TANDIS QUE LE
%  CHAMEAU (13) DOIT SE DECRIRE : « raƙumi mai tozo biyu », le chameau
%  a deux bosses. C'est l'inverse exact du francais, qui appelle
%  chameau la bete a deux bosses et doit specifier pour l'autre. Ici
%  la bete a une bosse EST le chameau, puisque c'est la seule qui
%  passe le desert — et c'est le meme raƙumi qui donnait la girafe et
%  la vague au tableau 10, « raƙumin dawa » et « raƙumin ruwa ». Trois
%  fois la meme bete pour mesurer le monde.
%
%  DOUZIEME REFUS DE LA COLONNE, ET NEUVIEME DU RELEVE SUR CETTE
%  PLANCHE : LE THEATRE DE MARIONNETTES (41). Le javanais avait
%  « wayang » et ne l'a pas ecrit ; le persan avait « خیمه‌شب‌بازی » et
%  ne l'a pas ecrit ; le haoussa a « tashe », la mascarade des nuits
%  de Ramadan, ou les enfants vont de maison en maison costumes et
%  jouent des saynetes pour un don. C'est ce que la langue a de plus
%  proche d'un theatre portatif, et ce n'est pas cela : la planche
%  grave un petit theatre a rideau rouge, decors peints, loges,
%  parterre et chef d'orchestre. On ecrit « gidan wasan ƴan tsana ».
%  TROIS COLONNES, TROIS MOTS TOUT PRETS, TROIS REFUS AU MEME RENVOI.
%
%  LE VIOLON (76) TIENT LA PROMESSE DU TABLEAU 11 : « bayolin », le v
%  refait en b, parce que le boko n'a pas de v. Et la meme lettre
%  refait la violette (94) en « bayolet », dans le meme tableau, a
%  vingt lignes de distance. Une contrainte d'alphabet appliquee deux
%  fois sans qu'on ait rien a decider.
%
%  LES QUATRE POINTS CARDINAUX (c) SONT LA, SUR LA ROSE DES VENTS DU
%  TABLEAU D'ASTRONOMIE, ET ILS CONFIRMENT LE TABLEAU 8 : Arewa,
%  Kudu, Gabas, Yamma, quatre mots haoussa. Le persan, au meme renvoi
%  du meme tableau, partageait les siens entre l'arabe et l'iranien.
%  Une langue qui a emprunte ses sept jours et ses douze mois n'a pas
%  emprunte une seule de ses quatre directions, et la planche
%  d'astronomie — c'est-a-dire ce qui S'ECRIT — est justement l'endroit
%  ou l'on aurait attendu l'emprunt.
%
%  ENFIN LE SIEGE DE L'ALINEA 5 EST HAOUSSA D'UN BOUT A L'AUTRE :
%  yaƙi la guerre, ganuwa l'enceinte — celle du tableau 11 —, sansani
%  le camp, hari l'assaut, ƙawanya le blocus, fatattaka la deroute,
%  ganima le butin, kama les prisonniers, sulke l'armure du tableau
%  10. Les cites haoussa se sont assiegees les unes les autres pendant
%  quatre siecles, et le grand-pere invalide raconte sa campagne dans
%  une langue qui n'a pas eu a chercher un mot.
% ===================================================================

%%K t16-apar-1 apar
\VUsaut{4.0mm}
\VUtitre{15.0pt}{60}{HOTO Na 16}
\VUsaut{2.0mm}
\VUpk{13.2pt}{40}{Babban shago. --- Kasuwar biki.}
\VUpk{13.2pt}{40}{Kayan wasa.}
\VUsaut{3.4mm}

%%K t16-01-1 p
1. --- Sabuwar shekara tana gabatowa. Manyan shaguna sun shirya
wuraren baje \VUgras{kayan wasansu} \textsuperscript{(1)} da kyautuka na sabuwar
shekara.

%%K t16-01-2 p
Na roƙi kakana ya kai ni kasuwar biki domin in ga wannan kyakkyawan
wurin baje koli, shi kuwa ya yarda.

%%K t16-02-1 p
2. --- Da farko muka tsaya gaban kyawawan garkuwoyin makamai, waɗanda
suke ɗauke da \VUgras{bindiga,} da \VUgras{hular soja,} da \VUgras{takobi
mai lanƙwasa,} da \VUgras{ɗamarar takobi} da \VUgras{alamun kafaɗa} biyu,
ko \VUgras{hular ƙarfe,} da \VUgras{sulke} \textsuperscript{(3)} da \VUgras{ƙaramar
bindiga} \textsuperscript{(4)}. Dukan wannan ya ba ni sha’awa fiye da
\VUgras{fitilun haske} \textsuperscript{(2)} ƙwarai.

%%K t16-tit-1 sub
\VUsaut{3.0mm}
\VUpk{10.2pt}{40}{Abubuwan kallo masu kyau.}
\VUsaut{2.6mm}

%%K t16-03-1 p
3. --- A wannan sashe akwai \VUgras{mai gadi} \textsuperscript{(5)} a cikin
\VUgras{rumfarsa} ; ya yi kama da wanda yake gadi gaban dukan gidan
namun daji na kwali. Kai, dabbobi nawa ne a can : \VUgras{aku}
\textsuperscript{(6)} a kan sandarsa, da \VUgras{karkanda} \textsuperscript{(7)} mai
ƙaho a kan hanci, da \VUgras{barewar ƙasar sanyi} \textsuperscript{(8)}, da
\VUgras{jimina} \textsuperscript{(9)}, da \VUgras{biri} \textsuperscript{(10)} yana
karya goro, da \VUgras{dila} \textsuperscript{(11)} yana ɗauke da kaza, da
\VUgras{kada} \textsuperscript{(12)} yana buɗe manyan muƙamuƙansa, da
\VUgras{raƙumi mai tozo biyu} \textsuperscript{(13)}, da \VUgras{karen farauta
mai sirara} \textsuperscript{(14)} mai kwarjini, da \VUgras{baƙar damisa}
\textsuperscript{(15)}, sa’an nan dabbobi biyu da ban sani ba waɗanda, bisa
ga faɗa, su ne \VUgras{ƙaramar dabba mai tsawo} \textsuperscript{(16)} da
\VUgras{kura} \textsuperscript{(17)}. Akwai kuma a can \VUgras{damisa}
\textsuperscript{(18)} da \VUgras{kerkeci} \textsuperscript{(21)} ; kusa ƙwarai akwai
kyawawan \VUgras{ƙananan kusu} \textsuperscript{(19)} da \VUgras{ƙananan ɓera}
\textsuperscript{(20)} na roba. A ƙarshe \VUgras{dorina}
\textsuperscript{(22)} da \VUgras{raƙumi} \textsuperscript{(23)} mai tozo suka gama
tarin. Da ma zan zauna a can na sa’o’i da yawa da ba a kusa da can
akwai kyakkyawan \VUgras{sansani} cike da \VUgras{sojojin dalma}
\textsuperscript{(29)} wanda ya ja hankalina ba.

%%K t16-tit-2 sub
\VUsaut{4.0mm}
\VUpk{10.2pt}{40}{Sojoji da yaƙi.}
\VUsaut{2.8mm}

%%K t16-04-1 p
4. --- Kakana, wanda yake \VUgras{naƙasasshe} \textsuperscript{(24)} kuma wanda
ya tilasta barin soja saboda \VUgras{rauni} da ya sa ya sami
\VUgras{lambar yabo ta soja,} yana son ya ba ni labarin \VUgras{yaƙe-yaƙen}
da ya shiga ƙwarai. Yana farin ciki yana bayyana mini motsi iri iri na
ƙaramar rundunar nan ƙanƙanuwa. Wata ƙungiyar sojoji na gaskiya da suke
ziyartar kasuwar bikin, wato \VUgras{soja mai gini} \textsuperscript{(25)}, da
\VUgras{mafarautan duwatsu} \textsuperscript{(26)} da kansa a rufe da
\VUgras{hula mai laushi,} da \VUgras{babban sajan} \textsuperscript{(27)} na sojojin
ƙasa da \VUgras{sajan igwa} \textsuperscript{(28)} mai \VUgras{alamar} zinare a
hannun riga, sun tsaya kusa da mu domin su saurari bayanan.

%%K t16-05-1 p
5. --- Kakana, cikin alfahari ƙwarai, tun da yake yana da masu
sauraro masu haske haka, sai ya ƙirƙiri cikakken shirin fada nan take.

%%K t16-05-2 p
Garin da aka yi wa ganuwa, tare da \VUgras{ganuwarsa} da
\VUgras{ramukansa,} \VUgras{rundunar abokan gaba} ta yi masa
\VUgras{ƙawanya,} wadda, bayan \VUgras{harbin igwa} mai tsawon lokaci, ta
kasance a shirye ta yi \VUgras{hari.} Amma \VUgras{waɗanda aka yi wa
ƙawanya,} da yunwa ta tilasta musu, sun yi ƙoƙarin \VUgras{fita da fada}
zuwa \VUgras{sansanin masu ƙawanyar.} Bayan sun fatattaki \VUgras{harin}
da \VUgras{fada} mai tsanani kewaye da \VUgras{sansanin,} masu ƙawanyar da
suka \VUgras{yi nasara} sun saki \VUgras{sojojin dokinsu} domin su kori
\VUgras{waɗanda aka ci da yaƙi.} Waɗannan sun \VUgras{ja da baya} suna
rufe \VUgras{filin fada} da \VUgras{matattu} da \VUgras{waɗanda suka ji
rauni.} \VUgras{Masu ɗaukar rauni} sun kai waɗannan na ƙarshe zuwa
\VUgras{motar asibiti} domin \VUgras{likitan tiyata} ya kula da su.

%%K t16-05-3 p
Duk da tsayin daka na jaruntaka na waɗansu \VUgras{bataliyoni} da suka
yi \VUgras{murabba’i,} \VUgras{shan kaye} ya cika, sauran rundunar sun
\VUgras{gudu,} suna barin \VUgras{kamammu} da yawa. Abokan gaba sun je su
cika \VUgras{nasararsu} suna mamaye garin. Wataƙila wannan zai zama
ƙarshen yaƙin kuma, bayan \VUgras{tsagaita wuta,} zai yiwu a sa hannu a
kan \VUgras{yarjejeniyar zaman lafiya.}

%%K t16-tit-3 sub
\VUsaut{4.4mm}
\VUpk{10.2pt}{40}{Kyautuka na sabuwar shekara.}
\VUsaut{2.8mm}

%%K t16-06-1 p
6. --- Matasan sojojin, waɗanda suka yi dariya suna sauraron labarin
tsohon sojan mai ban sha’awa, sun tambaye shi waɗansu bayanai. Game da
ni, na zagaya wurin sayarwa domin in kalli kyakkyawan \VUgras{bakan
harbi} \textsuperscript{(32)} da \VUgras{baka} \textsuperscript{(33)} tare da
\VUgras{kibiyarsa} \textsuperscript{(31)}. A can na sami wani tarin dabbobi :
\VUgras{tsuntsayen waƙa} \textsuperscript{(34)} biyu, wato \VUgras{tsuntsun waƙar
dare} na inji da \VUgras{tsuntsun kurmi} yana waƙa da dukan maƙogwaronsa,
da kuma jerin dabbobi masu ban mamaki : \VUgras{jemage}
\textsuperscript{(35)}, da \VUgras{mesa} \textsuperscript{(36)}, da \VUgras{kunkuru}
\textsuperscript{(37)}, da \VUgras{karen teku} \textsuperscript{(38)}, da
\VUgras{babban kifin teku} \textsuperscript{(39)} da \VUgras{kifi mai cin
mutane} \textsuperscript{(40)} na kwali.

%%K t16-07-1 p
7. --- Ban daɗe ina kallonsu ba, domin na hangi abin da nake mafarki :
kyakkyawan \VUgras{gidan wasan ƴan tsana} \textsuperscript{(41)} tare da
kyawawan \VUgras{kayan ado} da \VUgras{labule} ja mai ban sha’awa. A kan
\VUgras{fage,} ƴan wasa biyu sun yi kama da masu yin \VUgras{rawar gani}
a cikin \VUgras{wasan kwaikwayo} mai ban sha’awa ƙwarai, domin ƙananan
\VUgras{masu kallo} na kwali da aka ajiye a cikin ɗakuna ko zaune a kan
\VUgras{manyan kujeru} da a \VUgras{benen ƙasa} a bayan \VUgras{shugaban
ƙungiyar makaɗa,} sun yi ta tafi da ƙarfi.

%%K t16-07-2 p
Abin baƙin ciki, wannan gidan wasan yana da tsada ƙwarai.

%%K t16-08-1 p
8. --- Ban da haka, zaɓen kayan wasa yana da wuya, domin a cikin
tagogin an tara kowane iri kayan wasa : \VUgras{Arlekino}
\textsuperscript{(42)}, da \VUgras{Fulcinelo} \textsuperscript{(43)}, da
\VUgras{Fiyero} \textsuperscript{(44)}, da \VUgras{mujiyoyi} \textsuperscript{(45)} suna
kaɗa fikafikai, da \VUgras{baƙaƙen tsuntsu} \textsuperscript{(46)} masu busa, da
\VUgras{ƙananan karnuka} masu haushi, da \VUgras{injin waƙa}
\textsuperscript{(47)} da \VUgras{kayan wasa na haƙuri.} Ɗaya daga cikin
waɗannan kayan wasa ya ba ni daɗi ƙwarai : yana nuna \VUgras{fadar
jiragen ruwa} \textsuperscript{(48)} inda \VUgras{ƴan fashin teku} suke
kai wa wani \VUgras{jirgin ruwa} hari kusa da ƙasa da aka shuka
\VUgras{itatuwan kwakwa.}

%%K t16-noto-1 noto
\VUnotes{143.2mm}{%
\textit{(*) Duba hoto na 6.}%
}

%%K t16-09-1 p
9. --- Na iya kallon dukan waɗannan abubuwan mamaki cikin sauƙi ƙwarai,
domin mutane kaɗan ne kaɗai suke a shagon, da ƙyar waɗansu masu yawo,
da \VUgras{sojan doki} \textsuperscript{(49)} da \VUgras{mai karɓar kuɗi}
\textsuperscript{(50)} wanda ya gabatar da \VUgras{takardar biya} a babbar
\VUgras{ofishin kuɗi} \textsuperscript{(51)}.

%%K t16-10-1 p
10. --- Na tambayi kakana amfanin littattafan da aka ajiye a kan
kantuna a bayan \VUgras{ma’aikaciyar kuɗi} ; ya amsa mini cewa
\VUgras{littattafan lissafi} ne.

%%K t16-tit-4 sub
\VUsaut{3.6mm}
\VUpk{10.2pt}{40}{Kallo kewaye da shago.}
\VUsaut{2.8mm}

%%K t16-11-1 p
11. --- Muna barin sashen kayan wasa, sai muka nufi shagon sirdi domin
mu duba \VUgras{sirduna} \textsuperscript{(53)}, da \VUgras{ƙafafun sirdi}
\textsuperscript{(54)}, da \VUgras{linzamai} \textsuperscript{(55)}, da
\VUgras{bakunan linzami} \textsuperscript{(56)} da \VUgras{bulaloli.} Alhali
\VUgras{shugaban sashe} \textsuperscript{(57)} yake ba kakana waɗansu bayanai,
sai na je na kalli wurin baje \VUgras{kwanukan fenti} da
\VUgras{gilashi} \textsuperscript{(58)} inda kyawawan \VUgras{kwanonin salati} da
\VUgras{tukwanen shayi} \textsuperscript{(59)} masu laushi suke.

%%K t16-12-1 p
12. --- Domin mu je bene na farko, muka wuce a bayan tagar
\VUgras{matsattsun riguna} \textsuperscript{(60)}, kusa da wurin sayar da safa,
inda aka baje kyawawan \VUgras{wanduna} \textsuperscript{(63)} ƙwarai. A kusa,
wata \VUgras{ma’aikaciyar shago} \textsuperscript{(61)} tana sa wata \VUgras{mai
saye} \textsuperscript{(62)} ta zaɓi kyawawan \VUgras{saƙaƙƙun yadi}
\textsuperscript{(64)} na siliki.

%%K t16-12-2 p
Muna hawan matakala, na lura cewa an yi wa shagon ado da
\VUgras{fitilun hannu} \textsuperscript{(65)} na Japan, da \VUgras{laimu}
\textsuperscript{(66)} da manyan \VUgras{itatuwan dabino} \textsuperscript{(70)}.

%%K t16-13-1 p
13. --- A bene na farko, ba abin gani da yawa ; kayan gida kaɗai (*),
da \VUgras{akwatunan ƙarfe} \textsuperscript{(67)}, da \VUgras{kabad masu
aljihu} \textsuperscript{(68)} da \VUgras{abin hawan yara} \textsuperscript{(69)}. Amma
kallon shagon gaba ɗaya yana da \VUgras{daɗi} ƙwarai.

%%K t16-14-1 p
14. --- A ƙafafunmu, na ga kayan kiɗa : \VUgras{harfa}
\textsuperscript{(71)}, da \VUgras{gita} \textsuperscript{(72)}, da
\VUgras{mandolin} \textsuperscript{(73)}, da \VUgras{ƙananan busa}
\textsuperscript{(74)}, da \VUgras{sarewa masu ƙara} \textsuperscript{(75)}, da
bayolin tare da \VUgras{sandarsu} \textsuperscript{(76)} har ma da \VUgras{babbar
ganga} \textsuperscript{(77)}. Ɗan kaɗan can nesa, a sashen takarda, wani
\VUgras{ɗalibi} \textsuperscript{(78)} yana ciniki da \VUgras{ma’aikacin shago}
\textsuperscript{(79)} a kan jakar littattafai. Kusa da su, na rarrabe
kyakkyawan \VUgras{littafin haruffa} \textsuperscript{(80)}.

%%K t16-14-2 p
A tsakiyar shagon an kafa doguwar \VUgras{bishiyar Kirsimeti}
\textsuperscript{(81)} cike da kyandirori, da ƙananan abubuwa da kowane iri
kayan wasa : \VUgras{dawakin katako} \textsuperscript{(82)}, da \VUgras{ƴan
tsana} \textsuperscript{(83)}, da sauransu.

%%K t16-14-3 p
\VUgras{Wurin sayar da turare} \textsuperscript{(84)} tare da \VUgras{maganin
haƙori} da \VUgras{turare} iri iri ya baza ƙamshi mai kyau ƙwarai wanda
ya iso gare mu.

%%K t16-tit-5 sub
\VUsaut{3.4mm}
\VUpk{10.2pt}{40}{Mun sayi wani abu.}
\VUsaut{2.8mm}

%%K t16-15-1 p
15. --- Tun da kakana ya fara ganin lokacin ya yi tsawo ƙwarai, sai
muka sauka da babbar matakala. Ya ɗan tsaya gaban \VUgras{allon ilimin
taurari} \textsuperscript{(85)} wanda yake nuna, ya gaya mini,
\VUgras{taurari masu wutsiya} \textsuperscript{(a)}, da \VUgras{husufin wata da
na rana} \textsuperscript{(b)}, da fitowar iska tare da \VUgras{ɓangarori huɗu
na duniya} \textsuperscript{(c)} \VUgras{(Arewa, Kudu, Gabas da Yamma)}, da
taswirar \VUgras{rabin duniya} biyu \textsuperscript{(d)}, da
\VUgras{taurarin masu zagayawa} \textsuperscript{(e)} da \VUgras{tsarin rana}
\textsuperscript{(f)}. Ban gane komai a cikin dukan wannan ba, kuma
kayan aikin \VUgras{ɗan aiken kai kaya} \textsuperscript{(86)} mai alamun
zinare wanda ya wuce ya fi ba ni sha’awa, da kyawawan \VUgras{fankunan
hannu} \textsuperscript{(87)} da suke kusa da ni, kusa da \VUgras{jakunan
kuɗi} \textsuperscript{(88)} da \VUgras{jakunan takarda} \textsuperscript{(89)}.

%%K t16-16-1 p
16. --- Na kuma yaba wa \VUgras{gashin tsuntsu} \textsuperscript{(90)} a kan
wurin baje kayan, da \VUgras{ƙyalluran ɗamara} \textsuperscript{(91)}, da
kyawawan \VUgras{furannin ƙera} \textsuperscript{(94)} da aka tsara cikin
kwanduna da kwarjini. \VUgras{Furen bayolet,} da \VUgras{furen ƙamshi,}
da \VUgras{furen hiyasint} \textsuperscript{(95)}, da \VUgras{furen jeraniyo}
\textsuperscript{(96)}, da \VUgras{furen lili} \textsuperscript{(97)} da
\VUgras{furen nasturshin} \textsuperscript{(98)} an kwafe su da kyau ƙwarai.

%%K t16-tit-6 sub
\VUsaut{4.4mm}
\VUpk{10.2pt}{40}{Kyautar kakana.}
\VUsaut{2.8mm}

%%K t16-17-1 p
17. --- Na roƙi kakana ya saya mini ɗaya daga cikin \VUgras{buroshin
haƙori} \textsuperscript{(92)} da suke cikin akwati kusa da \VUgras{allurar
gashi} \textsuperscript{(93)}. Ya yarda ni kuwa na je da kaina na biya
sayena a ofishin kuɗi, kusa da inda ake shirya babban \VUgras{kaya}
\textsuperscript{(100)} domin wani \VUgras{abokin ciniki}
\textsuperscript{(99)}.

%%K t16-18-1 p
18. --- Ba zato ba tsammani wani ɗan hargitsi ya tashi tsakanin
mutanen da suka tsaya kusa da \VUgras{kayan tagulla}
\textsuperscript{(101)} na fasaha. An kama \VUgras{ɓarawo}
\textsuperscript{(102)} yanzu, cikin laifinsa, da \VUgras{mai duba}
\textsuperscript{(103)}, alhali yake ƙoƙarin yi wa wani sata. An kula a kawo
ɗan sanda domin ya \VUgras{kama} shi, kuma daidai lokacin da muka fita,
ana kai mai laifin zuwa kurkuku.

\VUsaut{6.0mm}
\VUfilet{22mm}
